\documentclass[10pt]{article}
\usepackage[margin=0.7in]{geometry}
\usepackage{amsmath,amssymb,amsthm}
\setlength{\parskip}{4pt}
\setlength{\parindent}{0pt}

\begin{document}
\fontsize{8}{10}\selectfont

\section*{Midterm Study Notes}

\subsection*{Section 2.2 \& 2.3 Topics}

\textbf{1. Definition of \(\mathbf{a \mid b}\).}\\
For integers \( a \) and \( b \), \( a \mid b\) if there exists an integer \( k \) such that \( b = ak \). Example: \(4 \mid 20\) because \(20 = 4 \cdot 5\), but \(4 \nmid 18\).

\textbf{2. Simple Properties of Division.}\\
Using \( a \mid b \iff b = a k\):
\begin{itemize}
\item \emph{Transitivity}: If \(a \mid b\) and \(b \mid c\), then \(a \mid c\).
\item \emph{Additivity}: If \(a \mid b\) and \(a \mid c\), then \(a \mid (b + c)\).
\item \emph{Cancellation}: If \(a \mid b\) and \(a \mid c\), then \(a \mid (mb + nc)\) for any integers \(m,n\).
\end{itemize}

\textbf{3. Division Algorithm.}\\
For \(a\) and \(b \neq 0\), there exist unique \(q,r\) with \(0 \le r < |b|\) such that
\[
a = bq + r.
\]
Example: \(23 = 5 \cdot 4 + 3\).

\textbf{4. Euclidean Algorithm (Finding \(\gcd(a,b)\)).}\\
Repeatedly apply the Division Algorithm:
\[
\gcd(a,b) = \gcd(b,r),
\]
until remainder \(=0\). Example: \(\gcd(48,18)\): 
\[
48 = 18 \cdot 2 + 12,\quad 18 = 12 \cdot 1 + 6,\quad 12 = 6 \cdot 2 + 0 \implies \gcd(48,18)=6.
\]

\textbf{5. Extended Euclidean Algorithm (Solving \(ax + by = \gcd(a,b)\)).}\\
Keep track of remainders to express \(\gcd(a,b)\) as \(ax + by\). Example: \(\gcd(48,18)=6\).
\[
48 = 18\cdot 2 + 12 \quad \Rightarrow \quad 12 = 48 - 2 \cdot 18,
\]
\[
18 = 12\cdot 1 + 6 \quad \Rightarrow \quad 6 = 18 - 12 = 18 - (48 - 2 \cdot 18) = 3 \cdot 18 - 48.
\]
Hence \(6=(-1)\cdot48 + 3\cdot18\), so \(x=-1,y=3\).

\textbf{6. Finding All Divisors of an Integer.}\\
Example: For \(36\), the positive divisors are \(\{1,2,3,4,6,9,12,18,36\}\). Negatives are their negatives.

\textbf{7. Fundamental Theorem of Arithmetic (Prime Factorization).}\\
Every positive integer has a unique prime factorization. Example: \(84 = 2^2 \times 3 \times 7\).

\textbf{8. Definition of \(\mathbf{a \equiv b \pmod{m}}\).}\\
\(a \equiv b \pmod{m}\) if \(m \mid (a - b)\). Equivalently, \(a\) and \(b\) have the same remainder mod \(m\). Example: \(14 \equiv 2 \pmod{4}\).

\textbf{9. Simple Properties of Congruences.}\\
If \(a \equiv b \pmod{m}\) and \(c \equiv d \pmod{m}\), then:
\begin{itemize}
\item \(a + c \equiv b + d \pmod{m}\).
\item \(ac \equiv bd \pmod{m}\).
\item \(-a \equiv -b \pmod{m}\).
\end{itemize}

\textbf{10. Modular Arithmetic in \(\mathbf{Z}/m\mathbf{Z}\).}\\
\(\mathbb{Z}/m\mathbb{Z} = \{ [0],[1],\ldots,[m-1] \}\). Operations:
\[
[a] + [b] = [a+b], \quad -[a] = [-a], \quad [a]\cdot[b] = [ab].
\]
Example: In \(\mathbb{Z}/5\mathbb{Z}\), \([2] + [4] = [6] = [1]\).

\textbf{11. Inverses in \(\mathbf{Z}/m\mathbf{Z}\) Using Extended Euclidean Algorithm.}\\
If \(\gcd(a,m)=1\), then \([a]\) has a multiplicative inverse in \(\mathbb{Z}/m\mathbb{Z}\). Example: In \(\mathbb{Z}/7\mathbb{Z}\), to find \([3]^{-1}\), solve \(3x+7y=1\). One solution is \(x=5\), so \([3]^{-1}=[5]\).

\subsection*{Section 3.1 Topics}

\textbf{1. Definition of \(\operatorname{Sym}(X)\).}\\
\(\operatorname{Sym}(X)\) is the set of all bijective functions \(f: X \to X\). It forms a group under composition. Example: if \(X = \{1,2,3\}\), any permutation (rearrangement) is in \(\operatorname{Sym}(X)\).

\textbf{2. Notations for Permutations in \(S_n\).}\\
\emph{Two-row notation}:
\[
\sigma = \begin{pmatrix}1 & 2 & 3 & \dots & n \\ \sigma(1) & \sigma(2) & \sigma(3) & \dots & \sigma(n)\end{pmatrix}.
\]
\emph{Cycle notation}: \((a_1\,a_2 \dots a_r)\) means \(\sigma(a_1)=a_2, \sigma(a_2)=a_3,\dots,\sigma(a_r)=a_1\).
Example: 
\[
\sigma=\begin{pmatrix}1 & 2 & 3 & 4 \\ 3 & 1 & 4 & 2\end{pmatrix} 
= (1\,3\,4\,2).
\]

\textbf{3. Composing Two Permutations.}\\
\(\tau \circ \sigma\) means apply \(\sigma\) first, then \(\tau\). Example:
\[
\sigma=\begin{pmatrix}1 & 2 & 3 \\ 2 & 3 & 1\end{pmatrix},
\quad
\tau=\begin{pmatrix}1 & 2 & 3 \\ 3 & 1 & 2\end{pmatrix}.
\]
Then \(\tau(\sigma(1)) = \tau(2)=1,\;\tau(\sigma(2))=\tau(3)=2,\;\tau(\sigma(3))=\tau(1)=3\). So \(\tau\circ\sigma = (1\,2\,3)\).

\textbf{4. More Complicated Permutation Compositions.}\\
\emph{Cycle Inverse}: \((a_1\,a_2\,\dots a_r)^{-1} = (a_r\, a_{r-1}\,\dots\,a_1)\).\\
\emph{Conjugation}: \(\tau(a_1\,a_2\,\dots a_r)\tau^{-1} = (\tau(a_1)\,\tau(a_2)\,\dots\,\tau(a_r))\).
Example: \((1\,2\,3)^{-1}=(3\,2\,1)\).

\subsection*{Section 3.2, 3.3, 4.1 Topics}

\textbf{1. Axioms of Semigroups, Monoids, and Groups.}\\
\begin{itemize}
\item \emph{Semigroup}: A set \(S\) with an associative binary operation.
\item \emph{Monoid}: A semigroup with identity \(e\) where \(e \cdot a=a \cdot e=a\).
\item \emph{Group}: A monoid in which every element has an inverse.
\end{itemize}

\textbf{2. Proving Simple Properties from Axioms.}\\
Example: from associativity and inverses, we get \emph{cancellation}: if \(ab=ac\), then \(b=c\). In \((\mathbb{Z},+)\), identity is \(0\), inverse of \(a\) is \(-a\).

\textbf{4. The Dihedral Group \(\mathbf{D_n}\).}\\
Group of symmetries of a regular \(n\)-gon:
\[
D_n = \{1, r, r^2,\dots,r^{n-1}, s, sr, sr^2,\dots,sr^{n-1}\},
\quad r^n=1,\; s^2=1,\; srs=r^{-1}.
\]
Interpretation: \(r\) is rotation by \(360^\circ/n\); \(s\) a reflection. Example: \(D_4\) (square) has 8 elements: \(\{1,r,r^2,r^3,s,sr,sr^2,sr^3\}\).

\textbf{5. Definition of a Subgroup.}\\
\(H \subseteq G\) is a subgroup if:
\begin{itemize}
\item \(H\neq \emptyset\),
\item If \(a,b \in H\), then \(ab \in H\),
\item If \(a\in H\), then \(a^{-1}\in H\).
\end{itemize}

\textbf{6. Using the Subgroup Criterion to Find All Subgroups.}\\
Example: In \(\mathbb{Z}_{12}\) under addition, subgroups are \(\langle d\rangle\) for each \(d\) dividing 12. E.g. \(\langle4\rangle=\{0,4,8\}\).

\textbf{7. Order of an Element.}\\
The order of \(a\) is the smallest positive integer \(n\) such that \(a^n=e\). Also the size of the cyclic subgroup generated by \(a\). Example: In \((\mathbb{Z}/12\mathbb{Z},+)\), \([3]\) has order 4 because \(3 \cdot 4=12\equiv0\pmod{12}\).

\textbf{8. Left Coset \(\mathbf{aH}\).}\\
If \(H\) is a subgroup of \(G\) and \(a\in G\), then \(aH=\{ah : h\in H\}\). Example: In \(\mathbb{Z}_{12}\), \(H=\{0,4,8\}\), then \([3]+H=\{[3],[7],[11]\}\).

\textbf{9. Lagrange’s Theorem.}\\
If \(H\) is a finite subgroup of \(G\), \(|H|\) divides \(|G|\). Thus if \(|G|=12\), subgroups can only have sizes dividing 12: \(1,2,3,4,6,12\).

\end{document}
